\documentclass{article}
\usepackage[utf8]{inputenc}
\usepackage[T1]{fontenc}
\usepackage{geometry}
\geometry{
    a4paper,
    left=2.5cm,
    right=2.5cm,
    top=2.5cm,
    bottom=2.5cm
}
\usepackage{graphicx}
\usepackage{float}
\usepackage{hyperref}
\usepackage[english]{babel}
\usepackage{amsmath}
\usepackage{amssymb}
\usepackage{array}
\usepackage{booktabs}
\usepackage{caption}
\usepackage{subcaption}
\usepackage{enumitem}
\usepackage{fancyhdr}
\usepackage{titlesec}
\usepackage{url}
\usepackage{textcomp}
\usepackage{xcolor}
\usepackage{listings}
\usepackage{multirow}
\usepackage{longtable}
\usepackage{tabularx}
\usepackage{times}

% Set up page style
\pagestyle{fancy}
\fancyhf{}
\rfoot{\thepage}
\renewcommand{\headrulewidth}{0pt}

% Title formatting
\titleformat{\chapter}[display]{\normalfont\bfseries}{\chaptertitle\thechapter}{0pt}{\Huge}
\titlespacing*{\chapter}{0pt}{50pt}{40pt}

% Document info
\title{Template Document}
\author{Academic Support Agent}
\date{\today}

\begin{document}
\maketitle

\maketitle

\question Write down sets in order of containment.

\begin{sol}
  We pretend that equivalence classes are just numbers.
  \[
    \C \supset \R \supset \Q \supset \Z \supset \N \supset
    \P \not\supset (\GF[7] = \modulo[7])  \supset \{\nil\}
  \]% the empty line bellow is needed to avoid extra space.
  
\end{sol}

\question Find roots of $x^2- 8x = 9$.

\begin{sol}
  We proceed by factoring,
  \begin{align*}
    x^2- 8x - 9     & = 9-9         &  & \text{Subtract 9 on both sides.}         \\
    x^2- x + 9x - 9 & = 0           &  & \text{Breaking the middle term.}         \\
    (x - 1)(x + 9)  & = 0           &  & \text{Pulling out common } (x - 1).      \\
    x               & \in \{1, -9\} &  & f(x)g(x) = 0 \Ra f(x) = 0 \vee g(x) = 0.
  \end{align*}% the empty line bellow is needed to avoid extra space.
  
\end{sol}

\question Figure \ref{wheel} shows two cipher wheels. The left one is from Jeffrey Hoffstein, et al. \cite{hoffstein2008introduction} (pg. 3). Write a Python 3 program that uses it to encrypt: \texttt{FOUR SCORE AND SEVEN YEARS AGO}.

\begin{sol}
  \img<wheel>[0.26]{Cipher wheels.}{cipher.png, diagram.jpg}

  The Python program is given in listing \ref{cpr} and the encryption is given in table \ref{enc}.

  \lstinputlisting[language=Python, caption={Python 3 implementing figure \ref{wheel} left wheel.}, label=cpr]{code/prog.py}

  \tbl<enc>{Caesar cipher} {
    Plain Text  & FOUR & SCORE & AND & SEVEN & YEARS & AGO \\
    Cipher Text & KTZW & XHTWJ & FSI & XJAJS & DJFWX & FLT \\
  }
\end{sol}

% citations
\end{document}